\chapter{Wiring Diagram Explained (Repository Reference)}
\label{ch:wiring-diagram-explained}


\section{Diagram (Included)}
\label{sec:wiring-diagram}

% -------------------- FIGURE AS PROVIDED --------------------
\begin{figure}[H]
    \centering
    \begin{circuitikz}[scale=0.6, transform shape, font=\large]
        \ctikzset{
            label distance=-0pt,
            every label/.style={font=\scriptsize}
        }
        % =====================================================
        % ESP32 BLOCK
        % =====================================================

        \node[draw, minimum width=3cm, minimum height=12.5cm] (esp) at (-3,0) {ESP32};

        % =====================================================
        % LEFT SIDE (INPUTS)
        % =====================================================

        % I2C
        \node[left=0.7cm of esp.north west, yshift=-4.8cm] (i2c_sda) {SDA (21)};
        \node[left=0.7cm of esp.north west, yshift=-5.8cm] (i2c_scl) {SCL (22)};

        % Analog inputs
        \node[left=2.25cm of esp.center, yshift=5.5cm] (num1) {NUM1 (32)};
        \node[left=2.25cm of esp.center, yshift=4.5cm] (num2) {NUM2 (33)};
        \node[left=2.25cm of esp.center, yshift=3.5cm] (ain1) {AIN1 (34)};
        \node[left=2.25cm of esp.center, yshift=2.5cm] (ain2) {AIN2 (35)};

        % Connection lines with larger open circles
        \draw (num1) -- ++(-2,0) node[ocirc, scale=3.6]{};
        \draw (num2) -- ++(-2,0) node[ocirc, scale=3.6]{};
        \draw (ain1) -- ++(-2.1,0) node[ocirc, scale=3.6]{};
        \draw (ain2) -- ++(-2.1,0) node[ocirc, scale=3.6]{};

        % Indicators
        \node[left=1cm of esp.south west, yshift=4cm] (ind1) {IND1 (4)};
        \node[left=1cm of esp.south west, yshift=3cm] (ind2) {IND2 (5)};
        \node[left=1cm of esp.south west, yshift=2cm] (ind3) {IND3 (2)};
        \node[left=1cm of esp.south west, yshift=1cm] (ind4) {IND4 (0)};

        % =====================================================
        % RIGHT SIDE (OUTPUTS)
        % =====================================================

        % PWM outputs
        \node[right=0.5cm of esp.north east, yshift=-0.8cm] (pwm1) {PWM LX (18)};
        \node[right=0.5cm of esp.north east, yshift=-1.8cm] (pwm2) {PWM LY (19)};
        \node[right=0.5cm of esp.north east, yshift=-2.8cm] (pwm3) {PWM RX (23)};
        \node[right=0.5cm of esp.north east, yshift=-3.8cm] (pwm4) {PWM RY (25)};
        \node[right=0.5cm of esp.north east, yshift=-4.8cm] (pwm5) {PWM KL (26)};
        \node[right=0.5cm of esp.north east, yshift=-5.8cm] (pwm6) {PWM KR (27)};

        % Switch inputs
        \node[right=2cm of esp.center, yshift=-1.4cm] (sw1) {SW1 (16)};
        \node[right=2cm of esp.center, yshift=-2.4cm] (sw2) {SW2 (17)};
        \node[right=2cm of esp.center, yshift=-3.4cm] (sw3) {SW3 (13)};
        \node[right=2cm of esp.center, yshift=-4.4cm] (sw4) {SW4 (14)};
        \node[right=2cm of esp.center, yshift=-5.4cm] (sw5) {SW5 (12)};
        \node[right=2cm of esp.center, yshift=-6.4cm] (sw6) {SW6 (15)};



        % =====================================================
        % I2C SENSOR
        % =====================================================

        \node[draw, minimum width=1.5cm, minimum height=0.8cm] (sensor) at (-14,2) {I2C};

        \draw (sensor.east) -- ++(0.3,0) |- (i2c_sda);
        \draw (sensor.east) -- ++(0.3,0) |- (i2c_scl);

        % Pull-up starting points
        \coordinate (sda_r) at ($(i2c_sda)+(-3,0)$);
        \coordinate (scl_r) at ($(i2c_scl)+(-5,0)$);

        % Resistors
        \draw (sda_r) to[R,l={\scriptsize $4.7k\Omega$}, label distance=-1pt] ++(0,1.5) coordinate (vbus1);
        \draw (scl_r) to[R,l={\scriptsize $4.7k\Omega$}, label distance=-1pt] ++(0,2.5) coordinate (vbus2);

        % Join Vcc rail
        \draw (vbus1) -- ++(0,0.5) coordinate (vccbus);
        \draw (vbus2) |- (vccbus);

        % Vcc node
        \draw (vccbus) -- ++(0,0.6) node[VCC33]{\scriptsize 3.3V};

        % =====================================================
        % BATTERY VOLTAGE MONITOR (ADC INPUT)
        % =====================================================

        % Local reference for the circuit
        \coordinate (batbase) at ($(ain1)+(9,-14)$);

        % Upper resistor
        \draw ($(batbase)+(0,1.5)$)
        to[R,l_=\scriptsize $100k$]
        (batbase);

        % Lower resistor
        \draw (batbase)
        to[R,l_=\scriptsize $33k$]
        ($(batbase)+(0,-1.5)$);

        % Battery positive
        \draw ($(batbase)+(0,1.5)$) -- ++(0,1) node[battery1]{};

        % VBAT label
        \node[left] at ($(batbase)+(0,2.5)$) {\scriptsize VBAT};

        % Battery negative to ground
        \draw ($(batbase)+(0,-1.5)$) -- ++(0,-1) node[bigground]{};

        % ADC connection terminal
        \draw (batbase) -- ++(2,0)
        node[ocirc, scale=3.6, label=right:{AIN2}]{};

        % =====================================================
        % LM35 TEMPERATURE SENSOR (AIN2)
        % =====================================================

        % Local reference for the LM35 block
        \coordinate (lm35base) at ($(ain2)+(3.5,-12)$);

        % Sensor body
        \node[draw, minimum width=1.8cm, minimum height=1.2cm] (lm35) at (lm35base) {LM35};

        % VCC pin
        \draw (lm35.north) -- ++(0,0.8) node[VCC33]{3.3V};

        % GND pin
        \draw (lm35.south) -- ++(0,-0.8) node[bigground]{};

        % Output pin to ESP32 terminal
        \draw (lm35.east) -- ++(1.5,0)
        node[ocirc, scale=3.6, label=right:{AIN1}]{};

        % =====================================================
        % LDR SENSOR INPUT (NUM2)
        % =====================================================

        \coordinate (ldrbase) at ($(num2)+(-7,-14)$);

        \begin{scope}[yshift=-7cm]

            % LDR (top part)
            \draw ($(ldrbase)+(0,2.5)$)
            to[photoresistor,l_=\normalsize LDR]
            (ldrbase);

            % Fixed resistor
            \draw (ldrbase)
            to[R,l_=\normalsize $10k$]
            ($(ldrbase)+(0,-1.5)$);

            % VCC
            \draw ($(ldrbase)+(0,1.5)$) -- ++(0,1) node[VCC33]{3.3V};

            % GND
            \draw ($(ldrbase)+(0,-1.5)$) -- ++(0,-1) node[bigground]{};

            % Connection terminal
            \draw (ldrbase) -- ++(2,0)
            node[ocirc, scale=3.6, label=right:{NUM2}]{};

        \end{scope}
        % =====================================================
        % POTENTIOMETER INPUT (NUM1)
        % =====================================================

        % Local reference for the potentiometer block
        \coordinate (potbase) at ($(num1)+(-2,-15.5)$);

        % Vertical potentiometer
        \draw ($(potbase)+(0,1.5)$)
        to[pR,l_=\scriptsize $10k$, invert]
        ($(potbase)+(0,-1.5)$);

        % Top terminal to VCC
        \draw ($(potbase)+(0,0.5)$) -- ++(0,1) node[VCC33]{3.3V};

        % Bottom terminal to GND
        \draw ($(potbase)+(0,-0.5)$) -- ++(0,-1) node[bigground]{};

        % Wiper arrow horizontal line to connection terminal
        \draw (potbase) -- ++(2,0)
        node[ocirc, scale=3.6, label=right:{NUM1}]{};



        % =====================================================
        % LED OUTPUTS (COMMON GROUND BUS)
        % =====================================================

        % Ground bus
        \draw (5.8,-3) -- (7,-3) node[bigground]{};

        % LED circuits
        \draw (sw1) to[R,l^=$330\Omega$, label distance=-3pt] ++(4.5,0)
        to[led] ++(1.2,0) coordinate (led1) -- ++(0,-1);

        \draw (sw2) to[R,l^=$330\Omega$, label distance=-3pt] ++(4.5,0)
        to[led] ++(1.2,0) coordinate (led2) -- ++(0,-1);

        \draw (sw3) to[R,l^=$330\Omega$, label distance=-3pt] ++(4.5,0)
        to[led] ++(1.2,0) coordinate (led3) -- ++(0,-1);

        \draw (sw4) to[R,l^=$330\Omega$, label distance=-3pt] ++(4.5,0)
        to[led] ++(1.2,0) coordinate (led4) -- ++(0,-1);

        \draw (sw5) to[R,l^=$330\Omega$, label distance=-3pt] ++(4.5,0)
        to[led] ++(1.2,0) coordinate (led5) -- ++(0,-1);

        \draw (sw6) to[R,l^=$330\Omega$, label distance=-3pt] ++(4.5,0)
        to[led] ++(1.2,0) coordinate (led6) -- ++(0,0);



        % =====================================================
        % SWITCH INPUTS (WITH PULL-UP AND COMMON VCC BUS)
        % =====================================================

        % -------- VCC BUS ----------
        \draw (-7.78,-0.29) -- (-14.5,-0.29) node[VCC33]{3.3V};

        % -----------------------------------------------------
        % IND1
        % -----------------------------------------------------
        \coordinate (ind1b) at ($(ind1)+(-1.2,0)$);
        \draw (ind1) -- (ind1b);

        \draw (ind1b) -- ++(0,0.0)
        to[R,l=$10k\Omega$] ++(0,2)
        coordinate (ind1r);

        \coordinate (sw1) at ($(ind1b)+(-1.5,0)$);
        \draw (sw1) node[bigground, line width=0.1pt]{}
        to[push button,l=\small BTN1] (ind1b);


        % -----------------------------------------------------
        % IND2
        % -----------------------------------------------------
        \coordinate (ind2b) at ($(ind2)+(-3.0,0)$);
        \draw (ind2) -- (ind2b);

        \draw (ind2b) -- ++(0,0.7)
        to[R,l=$10k\Omega$] ++(0,2.3)
        coordinate (ind2r);

        \coordinate (sw2) at ($(ind2b)+(-1.5,0)$);
        \draw (sw2) node[bigground, line width=0.1pt]{}
        to[push button,l=\small BTN2] (ind2b);


        % -----------------------------------------------------
        % IND3
        % -----------------------------------------------------
        \coordinate (ind3b) at ($(ind3)+(-4.8,0)$);
        \draw (ind3) -- (ind3b);

        \draw (ind3b) -- ++(0,1.4)
        to[R,l=$10k\Omega$] ++(0,2.55)
        coordinate (ind3r);

        \coordinate (sw3) at ($(ind3b)+(-1.5,0)$);
        \draw (sw3) node[bigground, line width=0.1pt]{}
        to[push button,l=\small BTN3] (ind3b);


        % -----------------------------------------------------
        % IND4
        % -----------------------------------------------------
        \coordinate (ind4b) at ($(ind4)+(-6.6,0)$);
        \draw (ind4) -- (ind4b);

        \draw (ind4b) -- ++(0,2.1)
        to[R,l=$10k\Omega$] ++(0,2.9)
        coordinate (ind4r);

        \coordinate (sw4) at ($(ind4b)+(-1.5,0)$);
        \draw (sw4) node[bigground, line width=0.1pt]{}
        to[push button,l=\small BTN4] (ind4b);
        % =====================================================
        % SERVOS
        % =====================================================
        \node[draw, minimum width=1.4cm, minimum height=1cm, anchor=west]
        (servo1) at ($(pwm1)+(5,0)$) {Servo};

        \node[draw, minimum width=1.4cm, minimum height=1cm, anchor=west]
        (servo3) at ($(pwm3)+(3.6,0)$) {Servo};

        % PWM signal
        \draw (pwm1) -- (servo1.west);
        \draw (pwm3) -- (servo3.west);

        % Power connections
        \draw (servo1.north) -- ++(0,0.2) node[VCC5]{5V};
        \draw (servo3.north) -- ++(0,0.2) node[VCC5]{5V};

        \draw (servo1.south) -- ++(0,-0.2) node[bigground]{};
        \draw (servo3.south) -- ++(0,-0.2) node[bigground]{};

        % Remaining PWM
        \draw (pwm5) -- ++(2.5,0) node[ocirc, scale=3.6]{};
        \draw (pwm6) -- ++(2.5,0) node[ocirc, scale=3.6]{};

        \draw (pwm2) -- ++(2.5,0) node[ocirc, scale=3.6]{};
        \draw (pwm4) -- ++(2.5,0) node[ocirc, scale=3.6]{};



    \end{circuitikz}
\caption{Suggested ESP32 pin assignment (part 1).}
\label{fig:full-rc-wiring-a}
\end{figure}

\clearpage % forces next floats to appear after a page break

% --- Figure 2: bottom part (your snippet) ---
\begin{figure}[H]
\centering
\begin{circuitikz}[scale=0.6, transform shape, font=\large]

        % =====================================================
        % PWM MOTOR CONTROL (PWM KL 26)
        % =====================================================

        % Local reference
        \coordinate (motorbase) at (-9.5,-22);

        % Motor symbol
        \node[draw, circle, minimum size=0.9cm] (motor) at ($(motorbase)+(0,2.0)$) {\small M};

        % MOSFET with part label
        \node[nmos, anchor=D, scale=2.4, transform shape,
        label=right:{\small IRLZ44N}]
        (mos) at ($(motorbase)+(0,-0.5)$) {};

        % Wire from MOSFET drain to motor bottom
        \draw (mos.D) -- (motor.south);

        % Wire from motor top to VCC
        \draw (motor.north) -- ++(0,1.8) node[VCC12]{6V to 12V};

        % MOSFET source to ground
        \draw (mos.S) -- ++(0,-0.8) node[bigground]{};

        % Flyback diode
        \draw (mos.D) -- ++(2.0,0)
        to[Do,l=\small 1N5819]
        ($(motor.north)+(2.0,1)$) -- ($(motor.north)+(0,1)$);

        % Gate resistor
        \draw ($(mos.G)+(-4,0)$)
        to[R,l=\small 220$\Omega$]
        (mos.G);

        % Pull-down resistor
        \draw ($(mos.G)+(-1,0)$)
        to[R,l=\small 10k]
        ($(mos.G)+(-1,-2.5)$)
        node[bigground]{};

        % PWM terminal
        \draw ($(mos.G)+(-4,0)$) -- ++(-0.5,0)
        node[ocirc, scale=3.6, label=left:{\small PWM KL (26)}]{};

        % =====================================================
        % PWM LED CONTROL (PWM KR 27)
        % =====================================================

        % Reference point
        \coordinate (pwmkr) at (-14.5,-18);

        % LED circuit
        \draw (pwmkr)
        to[R,l=\small 330$\Omega$] ++(2,0)
        to[led,l=\small LED] ++(1.5,0)
        -- ++(0,-1.2) node[bigground]{};

        % PWM control line
        \draw (pwmkr) -- ++(-1,0)
        node[ocirc, scale=3, label=left:{\small PWM KR (27)}]{};


        % =====================================================
        % H-BRIDGE MOTOR DRIVER (DRV8833) FOR PWM CONTROL
        % =====================================================

        \node[draw, thick, minimum width=3.8cm, minimum height=4cm] (drv) at (1,-22) {DRV8833};

        % -----------------------------------------------------
        % Motor A
        % -----------------------------------------------------

        \node[draw, circle, minimum size=0.9cm] (motorA) at (5,-23.0) {M};

        \draw (drv.east) ++(1.6,-0.4)
        node[left]{\small AOUT1}
        -| (motorA.north);

        \draw (drv.east) ++(1.6,-1.6)
        node[left]{\small AOUT2}
        -| (motorA.south);

        \node at (6.6,-23.0) {\small DC Motor A};

        % -----------------------------------------------------
        % Motor B
        % -----------------------------------------------------

        \node[draw, circle, minimum size=0.9cm] (motorB) at (5,-20.8) {M};

        \draw (drv.east) ++(1.6,1.8)
        node[left]{\small BOUT1}
        -| (motorB.north);

        \draw (drv.east) ++(1.6,0.6)
        node[left]{\small BOUT2}
        -| (motorB.south);

        \node at (6.6,-20.8) {\small DC Motor B};

        % -----------------------------------------------------
        % Motor supply
        % -----------------------------------------------------

        \draw ($(drv.north)+(0,0.5)$)
        node[below]{\small VM}
        -- ++(0,1)
        node[VCC12]{2.7V to 10.8V};

        % -----------------------------------------------------
        % Logic supply
        % -----------------------------------------------------

        \draw (drv.west) ++(-1.2,1.7)
        node[right]{\small VCC}
        -- ++(-1.5,0)
        node[vcc]{3.3V};

        % -----------------------------------------------------
        % Ground
        % -----------------------------------------------------

        \draw (drv.south)
        -- ++(0,-0.5)
        node[bigground,label={[yshift=-25]below:{\small GND}}]{};

        % -----------------------------------------------------
        % Motor A control
        % -----------------------------------------------------

        \draw (drv.west) ++(-1.2,-1)
        node[right]{\small AIN1}
        -- ++(-1.5,0)
        node[ocirc, scale=3.6,label=left:{\small PWM LY (19)}]{};

        \draw (drv.west) ++(-1.2,-1.6)
        node[right]{\small AIN2}
        -- ++(-1.5,0)
        node[ocirc, scale=3.6,label=left:{\small DIR (26)}]{};

        % -----------------------------------------------------
        % Motor B control
        % -----------------------------------------------------

        \draw (drv.west) ++(-1.2,0.6)
        node[right]{\small BIN1}
        -- ++(-1.5,0)
        node[ocirc, scale=3.6,label=left:{\small PWM RY (23)}]{};

        \draw (drv.west) ++(-1.2,-0.0)
        node[right]{\small BIN2}
        -- ++(-1.5,0)
        node[ocirc, scale=3.6,label=left:{\small DIR2 (27)}]{};

    \end{circuitikz}
    \caption{Suggested ESP32 pin assignment and peripheral connections for the \texttt{FULL\_RC\_MODE} configuration. The shown circuits are examples and can be modified according to the user's specific application requirements.}
    \label{fig:full-rc-wiring}
\end{figure}
% -------------------- END FIGURE --------------------



\section{Quick Summary}
\label{sec:wiring-quick-summary}

Figure~\ref{fig:full-rc-wiring} is a practical wiring reference for the \textbf{MICSBOL/ESP32\_BT\_Controller}
project. It shows:

\begin{itemize}
    \item \textbf{Inputs (left side):}
    \begin{itemize}
        \item I2C bus (SDA/SCL) for sensors or expansion
        \item Analog inputs (NUM1/NUM2/AIN1/AIN2) used as telemetry sources
        \item Digital indicator inputs (IND1--IND4) used to build a status mask
    \end{itemize}
    \item \textbf{Outputs (right side):}
    \begin{itemize}
        \item PWM outputs for servos, motors, LED brightness and buzzer
        \item Switch outputs (SW1--SW4 in direct mode, SW1--SW6 in MCP mode)
        \item Pulse/event outputs used for momentary actions (button events)
    \end{itemize}
\end{itemize}

\noindent In code terms:
\begin{itemize}
    \item \texttt{control.cpp} maps \texttt{rcStatePacket} fields to actions (steer, motors, pan, LED, buzzer, switches).
    \item \texttt{hal\_output.cpp} performs the real hardware writes (PWM via \texttt{ledcWrite}, GPIO via \texttt{digitalWrite},
    MCP outputs via \texttt{mcpWriteOutput}, pulses via \texttt{pulseStart}).
    \item \texttt{input\_hw.cpp} reads ADC and indicator pins (used for telemetry sources).
    \item \texttt{pins.h} defines the GPIO mapping (and it changes depending on the project mode).
\end{itemize}

\section{Straight Explanation of Each Group (What it does)}
\label{sec:wiring-components}

\subsection{I2C (SDA 21, SCL 22)}
Used to connect external I2C devices. In the repository, the same bus is also used by the
\texttt{FULL\_RC\_MODE\_MCP} configuration to talk to an MCP23017 expander (address \texttt{0x20} in \texttt{pins.h}).

\subsection{Analog telemetry inputs (NUM1/NUM2/AIN1/AIN2)}
These are ADC inputs read in \texttt{input\_hw.cpp} using \texttt{analogRead(...)}.
They are intended for telemetry sources such as battery voltage (through a divider), temperature sensors, light sensors,
potentiometers, etc.

\subsection{Digital indicator inputs (IND1--IND4)}
These are digital inputs (GPIO) that represent on/off statuses (limit switch, fault line, bumper, etc.).
In \texttt{input\_hw.cpp} they are combined into a bit mask so multiple statuses can be sent to the Android UI.

\subsection{PWM outputs (PWM LX/LY/RX/RY/KL/KR)}
PWM outputs are driven by \texttt{hal\_output.cpp} using \texttt{ledcWrite(...)} and are fed by values from
\texttt{control.cpp}.
Typical uses in the repository wiring concept:
\begin{itemize}
    \item servos (steering, camera pan)
    \item motor speed PWM
    \item LED brightness
    \item buzzer intensity
\end{itemize}

\subsection{Switch outputs (SW1--SW6)}
These are digital outputs controlled by switch bits received from the Android app (\texttt{control.cpp} calls
\texttt{halSetSwitch(index, state)}).
\begin{itemize}
    \item Direct mode provides 4 physical outputs (\texttt{PIN\_SW1}..\texttt{PIN\_SW4}).
    \item MCP mode provides 6 outputs via MCP23017 (\texttt{GPB0\_SW1}..\texttt{GPB5\_SW6}).
\end{itemize}

\subsection{Pulse/event outputs}
Android button events are handled by \texttt{controlHandleEvent(eventId)} in \texttt{control.cpp}, which triggers
momentary pulses via \texttt{halPulseSwitch(index)}. Pulses are implemented non-blocking in \texttt{pulse.cpp}.
These outputs are ideal for triggers (horn, latch, reset) rather than continuous states.

\section{Resume (One Paragraph)}
\label{sec:wiring-resume}

The wiring diagram groups the ESP32 connections exactly the way the repository code is organized:
\texttt{control.cpp} decides what each input means, \texttt{hal\_output.cpp} drives PWM and GPIO outputs,
\texttt{input\_hw.cpp} reads analog and indicator inputs for telemetry, and \texttt{pins.h} defines the GPIO mapping.
By keeping these responsibilities separated, you can adapt the hardware (change pin numbers, add I2C devices, switch to MCP23017)
without rewriting the overall control logic.
